%% ==================================================
%% Chapter: Matrix Product Operators and Density Operators
%% Introduces MPO and LPDO notation for the blueprint and records the
%% correspondence between local MPO data and the associated MPS view.
%% ==================================================

\chapter{Matrix Product Operators and Density Operators}\label{ch:mpdo}

Fix a physical dimension $d$ and a bond dimension $D$.
An MPO tensor carries two physical indices, one ket index and one bra index,
together with the virtual left and right indices.

\section{MPO tensors}

\begin{definition}[MPO tensor]\label{def:mpo_tensor}
    \lean{MPOTensor}
    \leanok
    An \emph{MPO tensor} is a family of matrices
    \[
        \{M^{ij}\}_{i,j=0}^{d-1}, \qquad M^{ij} \in \MN{D},
    \]
    indexed by a ket index $i$ and a bra index $j$. Diagrammatically,
    \[
        \TNMPOCell{M}{i}{j}
    \]
    The black node denotes the tensor $M$, the horizontal legs are virtual,
    the upper leg is the ket index $i$, and the lower leg is the bra
    index $j$.
\end{definition}

\begin{definition}[Doubled-index MPS view]\label{def:mpo_to_mps}
    \lean{MPOTensor.toMPSTensor}
    \leanok
    \uses{def:mpo_tensor}
    By combining the pair $(i,j)$ into a single physical index in
    $\{0,\ldots,d^2-1\}$, one obtains an MPS tensor with physical dimension
    $d^2$ and bond dimension $D$.
\end{definition}

\begin{definition}[Word evaluation]\label{def:mpo_eval_word}
    \lean{MPOTensor.evalWord}
    \leanok
    \uses{def:mpo_tensor}
    For words
    $u = (i_1,\ldots,i_N)$ and $v = (j_1,\ldots,j_N)$ in $\{0,\ldots,d{-}1\}^N$,
    the \emph{word evaluation} is
    \[
        M^{u,v} := M^{i_1 j_1} M^{i_2 j_2} \cdots M^{i_N j_N} \in \MN{D}.
    \]
    The empty pair of words evaluates to $\Id_D$, and word pairs of different
    lengths evaluate to $0$.
\end{definition}

\begin{definition}[MPO operator family]\label{def:mpo_family}
    \lean{MPOTensor.mpo}
    \leanok
    \uses{def:mpo_eval_word}
    The operator generated by $M$ on $N$ sites is the matrix
    $\rho^{(N)}(M)$ indexed by configurations
    $\sigma,\tau \in \{0,\ldots,d{-}1\}^N$ and given by
    \[
        \rho^{(N)}(M)_{\sigma,\tau} = \tr(M^{\sigma,\tau}).
    \]
\end{definition}

Diagrammatically,
\[
    \TNMPOChain{M}{\sigma_1}{\tau_1}{\sigma_N}{\tau_N}{N}
\]
each black node denotes the same local tensor $M$, and the matrix entry
$\rho^{(N)}(M)_{\sigma,\tau}$ is obtained by closing the outer virtual
legs cyclically and taking the resulting trace.

\begin{definition}[Hermitian MPO tensor]\label{def:mpo_hermitian}
    \lean{MPOTensor.IsHermitian}
    \leanok
    \uses{def:mpo_tensor}
    An MPO tensor is \emph{Hermitian} if
    \[
        M^{ij} = (M^{ji})^\dagger
        \qquad \text{for all } i,j \in \{0,\ldots,d{-}1\}.
    \]
\end{definition}

\section{Transfer maps}

\begin{definition}[MPO transfer map]\label{def:mpo_transfer_map}
    \lean{MPOTensor.transferMap}
    \leanok
    \uses{def:mpo_tensor}
    The \emph{transfer map} associated to $M$ is the linear map
    $\E_M : \MN{D} \to \MN{D}$ defined by
    \[
        \E_M(X) = \sum_{i,j=0}^{d-1} M^{ij} X (M^{ij})^\dagger.
    \]
\end{definition}

\begin{lemma}[Identification with the doubled-index MPS transfer map]\label{lem:mpo_transfer_eq_mps}
    \lean{MPOTensor.transferMap_eq_toMPSTensor}
    \leanok
    \uses{def:mpo_to_mps, def:mpo_transfer_map, def:transfer_map}
    The transfer map of an MPO tensor agrees with the transfer map of its
    doubled-index MPS view.
\end{lemma}

\begin{proof}
    \leanok
    \uses{def:mpo_to_mps, def:mpo_transfer_map, def:transfer_map}
    After identifying the single physical index of the doubled tensor with a
    pair $(i,j)$, the single sum defining the MPS transfer map is exactly the
    double sum defining $\E_M$.
\end{proof}

\begin{theorem}[Positivity of the MPO transfer map]\label{thm:mpo_transfer_pos}
    \lean{MPOTensor.transferMap_pos}
    \leanok
    \uses{def:mpo_transfer_map}
    If $X \ge 0$, then $\E_M(X) \ge 0$.
\end{theorem}

\begin{proof}
    \leanok
    \uses{lem:mpo_transfer_eq_mps}
    By Lemma~\ref{lem:mpo_transfer_eq_mps}, $\E_M$ is the transfer map of the
    doubled-index MPS tensor. The corresponding MPS transfer map is positive,
    so $\E_M(X) \ge 0$ whenever $X \ge 0$.
\end{proof}

\section{Zero correlation length}

\begin{definition}[Zero correlation length for MPO tensors]%
    \label{def:mpo_is_zcl}
    \lean{MPOTensor.IsZCL}
    \leanok
    \uses{def:mpo_transfer_map}
    An MPO tensor $M$ has \emph{zero correlation length} when its transfer map
    is idempotent:
    \[
        \E_M \circ \E_M = \E_M.
    \]
    This is the mixed-state analogue of the renormalization fixed-point
    condition for MPS tensors; see \cite[Definition~4.2]{Cirac2021Matrix}.
\end{definition}

\begin{theorem}[MPO ZCL iff doubled-index MPS RFP]%
    \label{thm:mpo_is_zcl_iff_to_mps_rfp}
    \lean{MPOTensor.isZCL_iff_toMPSTensor_isRFP}
    \leanok
    \uses{def:mpo_is_zcl, def:mpo_to_mps, def:is_rfp, lem:mpo_transfer_eq_mps}
    An MPO tensor has zero correlation length if and only if its doubled-index
    MPS view is a renormalization fixed point.
\end{theorem}

\begin{proof}
    \leanok
    By Lemma~\ref{lem:mpo_transfer_eq_mps}, the transfer map of $M$ agrees with
    the transfer map of its doubled-index MPS view. Thus
    $\E_M \circ \E_M = \E_M$ holds exactly when the transfer map of
    $M.\mathrm{toMPSTensor}$ is idempotent, which is the RFP condition.
\end{proof}

\section{MPDOs and LPDOs}

\begin{definition}[MPDO]\label{def:mpdo}
    \lean{MPOTensor.IsMPDO}
    \leanok
    \uses{def:mpo_family}
    An MPO tensor is an \emph{MPDO} if
    \[
        \rho^{(N)}(M) \ge 0
        \qquad \text{for every } N \in \N.
    \]
    This is the global positivity condition of \cite[\S4.1]{Cirac2017MPDO_AnnPhys}.
\end{definition}

\begin{definition}[LPDO]\label{def:lpdo}
    \lean{MPOTensor.IsLPDO}
    \leanok
    \uses{def:mpo_tensor}
    An MPO tensor is an \emph{LPDO} if there exist an ancilla dimension $d_K$,
    an inner bond dimension $D'$, an identification of bond indices
    $e : \{0, \ldots, D - 1\} \simeq \{0, \ldots, D' - 1\} \times \{0, \ldots, D' - 1\}$
    (equivalently, $D$ and $(D')^2$
    have the same cardinality), and matrices $A^{(i,k)} \in \MN{D'}$ such that
    \[
        M^{ij}
        = \Bigl(\sum_{k=0}^{d_K-1} A^{(i,k)} \otimes \overline{A^{(j,k)}}\Bigr)_{e,e}
        \qquad \text{for all } i,j,
    \]
    where $\otimes$ denotes the Kronecker product, $\overline{(\cdot)}$
    denotes entrywise complex conjugation, and $(\cdot)_{e,e}$ denotes
    reindexing along $e$ from $\{0, \ldots, D' - 1\} \times \{0, \ldots, D' - 1\}$
    back to $\{0, \ldots, D - 1\}$.
    This is the local purification form of \cite[\S4.3]{Cirac2017MPDO_AnnPhys}.
\end{definition}

\begin{theorem}[LPDOs are MPDOs]\label{thm:lpdo_is_mpdo}
    \lean{MPOTensor.IsLPDO.isMPDO}
    \leanok
    \uses{def:lpdo, def:mpdo}
    Every LPDO tensor is an MPDO.
\end{theorem}

\begin{proof}
    \leanok
    \uses{def:lpdo, def:mpdo}
    By the mixed-product property
    $(A \otimes B)(C \otimes D) = (AC) \otimes (BD)$, the $N$-site product
    decomposes as a sum of Kronecker products after transporting along the
    bond-space identification $e$:
    \[
        M^{\sigma_1 \tau_1} \cdots M^{\sigma_N \tau_N}
        = \bigl(\sum_{\kappa} P_\kappa \otimes \overline{Q_\kappa}\bigr)
          \big|_{e},
    \]
    where $P_\kappa = A^{(\sigma_1,\kappa_1)} \cdots A^{(\sigma_N,\kappa_N)}$
    and $Q_\kappa = A^{(\tau_1,\kappa_1)} \cdots A^{(\tau_N,\kappa_N)}$.
    Since the trace is invariant under the reindexing by~$e$,
    we apply $\tr(X \otimes Y) = \tr(X)\tr(Y)$ to obtain
    \[
        \rho^{(N)}(M)_{\sigma,\tau}
        = \sum_\kappa \tr(P_\kappa)\,\overline{\tr(Q_\kappa)}
        = \sum_\kappa \psi_\kappa(\sigma)\,\overline{\psi_\kappa(\tau)},
    \]
    where $\psi_\kappa(\sigma) = \tr(P_\kappa)$.
    Hence $\rho^{(N)}(M) = \sum_\kappa \ketbra{\psi_\kappa}{\psi_\kappa} \ge 0$.
\end{proof}

\section{Purification RFP}

\begin{definition}[Purifying MPS tensor]\label{def:purifying_mps_tensor}
    \lean{MPOTensor.purifyingMPSTensor}
    \leanok
    \uses{def:mpo_tensor}
    Given a purifying family
    $A^{(i,k)} \in \MN{D'}$ for $i \in \{0,\ldots,d{-}1\}$ and
    $k \in \{0,\ldots,d_K{-}1\}$, the \emph{purifying MPS tensor} is the
    MPS tensor with physical dimension $d \cdot d_K$ and bond dimension $D'$
    obtained by bundling the pair $(i,k)$ into a single index.
\end{definition}

\begin{definition}[Purification RFP]\label{def:is_prfp}
    \lean{MPOTensor.IsPRFP}
    \leanok
    \uses{def:lpdo, def:purifying_mps_tensor, def:is_rfp}
    An MPO tensor $M$ has a \emph{purification RFP} if it admits an LPDO
    witness whose purifying MPS tensor is a pure-state renormalization
    fixed point.
    This is \cite[Definition~4.3]{Cirac2021Matrix}.
\end{definition}

\begin{remark}[PRFP does not imply MPO ZCL]\label{rem:prfp_not_implies_zcl}
    \lean{MPOTensor.exists_isPRFP_not_isZCL}
    \leanok
    \uses{def:is_prfp, def:mpo_is_zcl}
    The current purification-RFP condition is strictly weaker than MPO zero
    correlation length: there exists an MPO tensor with purification RFP whose
    transfer map is not idempotent. Equivalently, the implication from the Lean
    predicate \texttt{IsPRFP} to the predicate \texttt{IsZCL} fails.
\end{remark}

\begin{theorem}[Purification RFP implies LPDO]\label{thm:prfp_is_lpdo}
    \lean{MPOTensor.IsPRFP.isLPDO}
    \leanok
    \uses{def:is_prfp, def:lpdo}
    A purification-RFP tensor is, in particular, an LPDO.
\end{theorem}

\begin{proof}
    \leanok
    The PRFP condition provides an LPDO witness by construction.
\end{proof}

\section{Pure-state recovery inside the MPO formalism}

\begin{definition}[MPDO renormalization fixed point]\label{def:is_rfp_mpdo}
    \lean{MPOTensor.IsRFP}
    \leanok
    \uses{def:mpo_transfer_map}
    An MPO tensor $M$ is an \emph{MPDO renormalization fixed point} if its
    transfer map is idempotent:
    \[
        \E_M \circ \E_M = \E_M.
    \]
\end{definition}

\begin{theorem}[MPDO RFP iff MPO ZCL]\label{thm:is_rfp_iff_is_zcl}
    \lean{MPOTensor.isRFP_iff_isZCL}
    \leanok
    \uses{def:is_rfp_mpdo, def:mpo_is_zcl}
    The provisional MPDO renormalization fixed-point condition is equivalent to
    the zero-correlation-length condition for MPO tensors.
\end{theorem}

\begin{proof}
    \leanok
    Both conditions assert $\E_M \circ \E_M = \E_M$; they are definitionally
    identical.
\end{proof}

\begin{theorem}[MPDO RFP iff doubled-index MPS RFP]\label{thm:is_rfp_iff_to_mps_rfp}
    \lean{MPOTensor.isRFP_iff_toMPSTensor_isRFP}
    \leanok
    \uses{def:is_rfp_mpdo, def:mpo_to_mps, def:is_rfp, def:mpo_is_zcl,
        thm:mpo_is_zcl_iff_to_mps_rfp}
    The provisional MPDO RFP condition is equivalent to the pure-state RFP
    condition for the doubled-index MPS tensor.
\end{theorem}

\begin{proof}
    \leanok
    \uses{thm:mpo_is_zcl_iff_to_mps_rfp}
    Unfold the MPDO RFP condition to its definitional equivalent, MPO ZCL,
    then apply Theorem~\ref{thm:mpo_is_zcl_iff_to_mps_rfp}.
\end{proof}

\begin{definition}[Diagonal pure-state embedding as an MPO]\label{def:to_mpo_tensor}
    \lean{MPSTensor.toMPOTensor}
    \leanok
    \uses{def:mps_tensor, def:mpo_tensor}
    An MPS tensor $A = \{A^i\}_{i=0}^{d-1}$ determines an MPO tensor by
    placing $A^i$ on the diagonal in the bra and ket indices:
    \[
        M^{ij} = \delta_{ij} A^i.
    \]
    Equivalently, $M^{ii} = A^i$ and $M^{ij} = 0$ for $i \ne j$.
\end{definition}

\begin{theorem}[Transfer map of the diagonal pure-state embedding]%
    \label{thm:to_mpo_transfer_map}
    \lean{MPSTensor.toMPOTensor_transferMap}
    \leanok
    \uses{def:to_mpo_tensor, def:mpo_transfer_map, def:transfer_map}
    The transfer map of the diagonal MPO associated to $A$ agrees with the
    original MPS transfer map:
    \[
        \E_{A^{\mathrm{MPO}}} = \E_A.
    \]
\end{theorem}

\begin{proof}
    \leanok
    \uses{def:to_mpo_tensor, def:mpo_transfer_map, def:transfer_map}
    In the double sum defining $\E_{A^{\mathrm{MPO}}}$, all off-diagonal terms
    vanish because $M^{ij} = 0$ for $i \ne j$, while the diagonal terms are
    exactly $A^i X (A^i)^\dagger$. Thus only the sum over $i$ remains, which
    is the defining formula for $\E_A$.
\end{proof}

\begin{theorem}[Pure-state recovery of the RFP condition]%
    \label{thm:to_mpo_rfp_iff_rfp}
    \lean{MPSTensor.toMPOTensor_isRFP_iff_isRFP}
    \leanok
    \uses{def:is_rfp_mpdo, def:is_rfp, thm:to_mpo_transfer_map}
    For an MPS tensor $A$, its diagonal MPO embedding is MPDO-RFP if and only
    if $A$ is an RFP:
    \[
        A^{\mathrm{MPO}} \text{ is MPDO-RFP} \iff A \text{ is RFP}.
    \]
\end{theorem}

\begin{proof}
    \leanok
    \uses{def:is_rfp_mpdo, def:is_rfp, thm:to_mpo_transfer_map}
    Both conditions assert idempotence of the corresponding transfer map.
    Theorem~\ref{thm:to_mpo_transfer_map} identifies those transfer maps, so
    the two predicates are equivalent.
\end{proof}

\begin{theorem}[Pure-state recovery of zero correlation length]%
    \label{thm:to_mpo_rfp_iff_zcl}
    \lean{MPSTensor.toMPOTensor_isRFP_iff_isZCL}
    \leanok
    \uses{thm:to_mpo_rfp_iff_rfp, thm:zcl_iff_idempotent_transfer}
    For an MPS tensor $A$, the diagonal MPO embedding is MPDO-RFP if and only
    if $A$ has zero correlation length.
\end{theorem}

\begin{proof}
    \leanok
    \uses{thm:to_mpo_rfp_iff_rfp, thm:zcl_iff_idempotent_transfer}
    Combine Theorem~\ref{thm:to_mpo_rfp_iff_rfp} with the characterization of
    zero correlation length by idempotence of the MPS transfer map
    (Theorem~\ref{thm:zcl_iff_idempotent_transfer}).
\end{proof}

\section{Simple local structure from SAL and ZCL}

For the simple MPDO case in Appendix~C.2 of \cite{Cirac2017MPDO_AnnPhys},
the strong-area-law hypothesis is used locally through the equality case of
strong subadditivity for a normalized three-site reduced state.

\begin{definition}[Local $\eta$-structure for a three-site reduced state]
    \label{def:mpdo_eta_structure}
    \lean{MPOTensor.EtaStructure}
    \leanok
    \uses{def:entropy_quantum_markov_decomposition}
    For a three-site density operator $\rho_{ABC}$, the local $\eta$-structure
    is the quantum Markov decomposition on the middle subsystem $B$ produced by
    equality in strong subadditivity. Concretely, after a unitary change of
    basis on $B$, one has a finite direct-sum decomposition
    \[
        B = \bigoplus_k \bigl(b_1^{(k)} \otimes b_2^{(k)}\bigr)
    \]
    such that
    \[
        \rho_{ABC} = \bigoplus_k \rho_{A b_1}^{(k)} \otimes \rho_{b_2 C}^{(k)}.
    \]
\end{definition}

\begin{theorem}[SAL implies local $\eta$-structure]
    \label{thm:sal_implies_eta_structure}
    \lean{MPOTensor.sal_implies_eta_structure}
    \leanok
    \uses{def:ssa_equality, def:mpdo_eta_structure}
    If a normalized three-site reduced state satisfies the equality case of
    strong subadditivity --- the local entropy form of the simple-MPDO strong
    area law --- then it admits the local $\eta$-structure of
    Definition~\ref{def:mpdo_eta_structure}.
\end{theorem}

\begin{proof}
    \leanok
    \uses{thm:entropy_ssa_equality_quantum_markov}
    This is exactly the forward implication of the Hayashi equality
    characterization: equality in strong subadditivity yields a quantum Markov
    decomposition on the middle subsystem.
\end{proof}

\begin{definition}[Auxiliary matrix predicates for the rank-one step]
    \label{def:matrix_trace_rankone_predicates}
    \lean{Matrix.TracePowersConstant}
    \lean{Matrix.HasRankOneFactorization}
    \lean{Matrix.PrimitiveTracePowersConstantImpliesRankOne}
    \leanok
    The auxiliary predicates record three matrix-theoretic ingredients used in
    the rank-one step of Appendix~C.2, Lemma~C.4: constant traces of all
    positive powers, existence of a rank-one factorization $T = a b^\top$, and
    the Perron--Frobenius implication from primitivity together with constant
    trace powers to such a factorization.
\end{definition}

\begin{theorem}[SAL and ZCL imply rank-one $T$ (scoped form)]
    \label{thm:sal_zcl_implies_rank_one_t}
    \lean{MPOTensor.sal_zcl_implies_rank_one_T}
    \leanok
    \uses{def:matrix_trace_rankone_predicates}
    Let $T$ be the nonnegative matrix extracted from the local $\eta$-data.
    Assume that $T$ is primitive, that $\tr(T) = 1$, and that every positive
    power of $T$ has the same trace as $T$. If one further supplies the single
    Perron--Frobenius input asserting that every primitive nonnegative matrix
    with constant trace powers is rank one, then
    \[
        T = a b^\top
    \]
    for some vectors $a$ and $b$, with normalization
    \[
        a \cdot b = 1.
    \]
\end{theorem}

\begin{proof}
    \leanok
    \uses{def:matrix_trace_rankone_predicates}
    The Perron--Frobenius input gives the factorization $T = a b^\top$.
    Taking traces and using $\tr(a b^\top) = a \cdot b$ converts the
    normalization $\tr(T) = 1$ into $a \cdot b = 1$.
\end{proof}

\section{Horizontal and Vertical Canonical Forms}

\begin{definition}[Horizontal canonical form for an MPO]\label{def:horizontal_cf_mpdo}
    \lean{MPOTensor.IsHorizontalCF}
    \notready
    \uses{def:mpo_tensor, def:tensor_from_blocks, def:injective, def:same_mpv2}
    An MPO tensor is in \emph{horizontal canonical form} if its doubled-index
    MPS tensor generates the same MPV family as a block-diagonal tensor
    $\bigoplus_k \mu_k A_k$ whose blocks are injective, left-canonical, and
    have nonzero weights $\mu_k$.
\end{definition}

\begin{definition}[Vertical canonical form for an MPO]\label{def:vertical_cf_mpdo}
    \lean{MPOTensor.IsVerticalCF}
    \notready
    \uses{def:mpo_tensor, def:bnt, def:same_mpv2, def:tensor_from_blocks}
    An MPO tensor is in \emph{vertical canonical form} if its diagonal MPS
    tensor $\{M^{ii}\}_{i=0}^{d-1}$ generates the same MPV family as a
    block-diagonal tensor obtained from a basis of normal tensors by repeating
    each block according to a multiplicity and weighting each copy by a
    positive scalar.
    This is the flattened form of the decomposition
    $\bigoplus_\alpha \mu_\alpha \otimes M_\alpha$
    from \cite[\S4.4]{Cirac2017MPDO_AnnPhys}.
\end{definition}

\begin{remark}[Finite-length span and selector criteria for the abstract biCF hypothesis]%
    \label{rem:word_tuple_span_top}
    \lean{MPSTensor.WordTupleSpanTop}
    \lean{MPSTensor.HasBiCF}
    \lean{MPSTensor.hasBiCF_of_wordTupleSpanTop}
    \lean{MPSTensor.HasBlockSelectorWords}
    \lean{MPSTensor.PropBlockInjective}
    \lean{MPSTensor.wordTupleSpanTop_of_propBlockInjective}
    \lean{MPSTensor.hasBiCF_of_propBlockInjective}
    \lean{MPOTensor.horizontalCFData_of_wordTupleSpanTop}
    \lean{MPOTensor.horizontalCFData_of_propBlockInjective}
    \leanok
    Let $(A_k)$ be a family of block tensors, and fix a length $L$.
    Write $A_k^w$ for the length-$L$ word evaluation of $A_k$ along a word $w$.
    The predicate \texttt{WordTupleSpanTop} says that the tuple-valued evaluations
    \[
        w \longmapsto (A_k^w)_k
    \]
    span the full product algebra of block matrices.
    The predicate \texttt{HasBiCF} records the corresponding finite-length
    separation property: there exists $L$ such that if
    \[
        \sum_k \tr(\Delta_k A_k^w) = 0
        \qquad \text{for every word } w \text{ of length } L,
    \]
    then every block coefficient $\Delta_k$ vanishes.
    The theorem \texttt{hasBiCF\_of\_wordTupleSpanTop} proves that the
    spanning condition implies this separation property, by nondegeneracy of the
    product trace pairing. The packaging theorem
    \texttt{horizontalCFData\_of\_wordTupleSpanTop} then inserts the resulting
    biCF witness into \texttt{HorizontalCFData} once injectivity,
    left-canonicality, and nonzero weights are given.

    The predicate \texttt{PropBlockInjective} isolates the finite-length content
    of Proposition~IV.3: after some common blocking length, each block is
    block-injective, and after a second finite blocking length one can form word
    combinations which equal the identity on one chosen block and vanish on all
    others. Concatenating an injective prefix with such a selector suffix yields
    the product-algebra span condition above. This is formalized by
    \texttt{wordTupleSpanTop\_of\_propBlockInjective}, and the resulting biCF
    witness is packaged into horizontal canonical-form data by
    \texttt{horizontalCFData\_of\_propBlockInjective}.
\end{remark}

\begin{lemma}[Blockwise equality from agreement on the MPV family]
    \label{lem:blockwise_insert_eq_of_mpv_agree}
    \lean{MPOTensor.blockwise_insert_eq_of_mpv_agree}
    \leanok
    \uses{def:tensor_from_blocks, def:injective}
    Let $\bigoplus_k \mu_k A_k$ be a block-diagonal tensor whose blocks are
    injective, left-canonical, weighted by nonzero scalars, and jointly
    block-injective in the sense that for some blocking length~$L$ the tuple of
    trace pairings against length-$L$ block words is faithful across all blocks
    simultaneously (the block-injective canonical form, biCF).
    If two physical operators induce the same first-site action on every MPV
    generated by this tensor, then their first-site insertions agree on each
    block $A_k$.
\end{lemma}

\begin{proof}
    \leanok
    \uses{def:tensor_from_blocks, def:injective}
    This is Lemma~L from \cite[Appendix~A]{Cirac2017MPDO_AnnPhys}.
    Expanding the hypothesis over a word of length~$L+1$ beginning with the
    fixed site and folding each block contribution into a trace pairing
    against $\mathrm{insertedTensor}\,W\,A_k$ yields, upon taking the
    $Y$\!-$Z$ difference, a tuple of block matrices whose trace pairings
    against every length-$L$ block word vanish. The biCF hypothesis then
    forces each block difference to be zero, and dividing by the nonzero
    weight $\mu_k^{L+1}$ concludes the blockwise equality of the two
    first-site insertions.
\end{proof}

\begin{theorem}[Horizontal canonical form implies vertical canonical form]
    \label{thm:vertical_cf_of_horizontal_cf}
    \notready
    \uses{def:mpdo, def:horizontal_cf_mpdo, def:vertical_cf_mpdo,
        lem:blockwise_insert_eq_of_mpv_agree}
    Every MPDO in horizontal canonical form is also in vertical canonical form.
\end{theorem}

\begin{proof}
    \uses{def:mpdo, def:horizontal_cf_mpdo, def:vertical_cf_mpdo,
        lem:blockwise_insert_eq_of_mpv_agree}
    Starting from the horizontal canonical-form decomposition, one uses the MPDO
    positivity condition to compare the first-site action of the doubled-index
    tensor on the diagonal sector. Lemma~\ref{lem:blockwise_insert_eq_of_mpv_agree}
    then shows that the diagonal tensor splits
    along the same normal blocks, with the coefficients rearranged into
    positive scalar weights.
    Flattening those positive diagonal coefficients produces the required
    vertical canonical-form decomposition.
\end{proof}

\section{Fusion isometries}

Following \cite[\S4.5]{Cirac2017MPDO_AnnPhys}, the current Lean development
formalizes the transfer-map content of the fusion-isometry picture.
For each blocked size $n \ge 1$, the doubled-index blocked tensor of an MPO
has a blocked transfer map $E_n$ acting on bond-space matrices, and the
support algebra is modeled by a subspace through which $E_n$ factors as a
retract.

\begin{definition}[Transfer-map fusion-isometry datum]%
    \label{def:mpdo_fusion_isometry}
    \lean{MPOTensor.FusionIsometryData}
    \leanok
    \uses{def:mpo_tensor}
    A \emph{transfer-map fusion-isometry datum} at blocked size $n$ for an
    MPO tensor consists of a subspace $\mathcal{A}_n \subseteq M_D(\mathbb{C})$
    together with linear maps
    $$T_n : M_D(\mathbb{C}) \to \mathcal{A}_n, \qquad
      S_n : \mathcal{A}_n \to M_D(\mathbb{C})$$
    such that
    $$T_n \circ S_n = \operatorname{id}_{\mathcal{A}_n},
      \qquad S_n \circ T_n = E_n,$$
    where $E_n$ denotes the blocked transfer map.
\end{definition}

\begin{definition}[Fusion-isometry RFP predicate]%
    \label{def:mpdo_rfp_via_fusion}
    \lean{MPOTensor.IsRFP_MPDO_via_fusion}
    \leanok
    \uses{def:mpdo_fusion_isometry, def:mpo_tensor}
    An MPO tensor $M$ satisfies the \emph{fusion-isometry formulation} of the
    renormalization fixed-point condition when for every blocked size $n \ge 1$
    there exists transfer-map fusion-isometry data for the blocked transfer
    map $E_n$.
\end{definition}

\begin{theorem}[Fusion-isometry data imply the provisional MPDO RFP condition]%
    \label{thm:mpdo_rfp_of_fusion}
    \lean{MPOTensor.isRFP_of_isRFP_MPDO_via_fusion}
    \leanok
    \uses{def:mpdo_rfp_via_fusion}
    The fusion-isometry formulation implies the current predicate
    \lean{MPOTensor.IsRFP}.
\end{theorem}
\begin{proof}\leanok
Apply the definition at blocked size $n = 1$.
If $E_1 = S_1 \circ T_1$ and $T_1 \circ S_1 = \operatorname{id}$, then
$$E_1^2 = S_1 T_1 S_1 T_1 = S_1 (T_1 S_1) T_1 = S_1 T_1 = E_1.$$
Since $E_1$ is the original transfer map, this is exactly
\lean{MPOTensor.IsRFP}.
\end{proof}

\begin{theorem}[Transfer-map fusion criterion for MPDO RFP]%
    \label{thm:mpdo_rfp_via_fusion_iff}
    \lean{MPOTensor.isRFP_MPDO_via_fusion_iff_isRFP}
    \leanok
    \uses{def:mpdo_rfp_via_fusion}
    The fusion-isometry formulation is equivalent to
    \lean{MPOTensor.IsRFP}.
\end{theorem}
\begin{proof}\leanok\uses{thm:mpdo_rfp_of_fusion}
The forward implication is Theorem~\ref{thm:mpdo_rfp_of_fusion}.
Conversely, if the transfer map $E$ is idempotent, then every positive blocked
transfer map equals $E$ and therefore factors through its range.
This range-factorization gives the required fusion-isometry data.
\end{proof}

\begin{corollary}[Pure-state recovery]%
    \label{cor:mpdo_fusion_pure_recovery}
    \lean{MPSTensor.toMPOTensor_isRFP_MPDO_via_fusion_iff_isRFP}
    \leanok
    \uses{thm:mpdo_rfp_via_fusion_iff}
    For the diagonal MPO associated to a pure MPS tensor, the fusion-isometry
    formulation reduces to the usual pure-state RFP condition.
\end{corollary}
\begin{proof}\leanok
Combine Theorem~\ref{thm:mpdo_rfp_via_fusion_iff} with the diagonal pure-state
recovery theorem for the provisional MPO predicate.
\end{proof}

\section{Algebra structure}

\begin{definition}[Algebra-structure data for an MPDO]%
    \label{def:mpdo_algebra_structure_data}
    \lean{MPOTensor.AlgebraStructureData}
    \leanok
    \uses{def:mpo_tensor}
    A family of support algebras $\{\mathcal{A}_n\}_{n \ge 0}$, one at every
    blocking size, equipped with blocking maps
    $m_n : \mathcal{A}_n \otimes \mathcal{A}_n \to \mathcal{A}_{2n}$ arising
    from horizontal concatenation of blocks and unital inclusion maps
    $\iota_n : \mathcal{A}_n \to \mathcal{A}_{n+1}$ arising from appending a
    single site, together with a coherence condition relating $m_n$ and
    $\iota_n$.
\end{definition}

\begin{definition}[Algebra-structure RFP predicate (provisional)]%
    \label{def:mpdo_rfp_via_algebra}
    \lean{MPOTensor.IsRFP_MPDO_via_algebra_scaffold}
    \uses{def:mpdo_algebra_structure_data, def:mpo_tensor}
    An MPO tensor $M$ satisfies the \emph{algebra-structure formulation} of
    the renormalization fixed-point condition when there exist
    algebra-structure data compatible with $M$, in the sense of
    \cite[Theorem~IV.13~(ii)]{Cirac2017MPDO_AnnPhys}. The compatibility
    relation between algebra-structure data and~$M$ is, in this provisional
    formulation, the trivial relation that holds for every pair;
    consequently the predicate above holds vacuously of every MPO tensor.
    The exact compatibility condition of
    \cite[Theorem~IV.13~(ii)]{Cirac2017MPDO_AnnPhys} is to be substituted
    in place of the trivial relation in subsequent work.
\end{definition}

\section{Commuting form and GSNNCH}

\begin{definition}[Commuting-form data for an MPDO]
    \label{def:mpdo_commuting_form_data}
    \lean{MPOTensor.CommutingFormData}
    \leanok
    \uses{def:mpo_tensor}
    For a fixed chain length $N \ge 2$, this consists of a positive
    semidefinite two-site operator $B$ such that its translated copies on the
    periodic chain commute pairwise.
\end{definition}

\begin{definition}[Commuting-form property]
    \label{def:mpdo_has_commuting_form}
    \lean{MPOTensor.HasCommutingForm}
    \leanok
    \uses{def:mpdo_commuting_form_data, def:mpo_tensor}
    An MPO tensor has the commuting-form property if for every $N \ge 2$ its
    $N$-site operator is a positive scalar multiple of the product of the
    translated copies of some commuting two-site operator $B$.
\end{definition}

\begin{definition}[GSNNCH]
    \label{def:mpdo_is_gsnnch}
    \lean{MPOTensor.IsGSNNCH}
    \leanok
    \uses{def:mpdo_has_commuting_form}
    An MPO tensor is a Gibbs state of a nearest-neighbor commuting Hamiltonian
    if every finite-chain operator admits the equivalent positive-operator
    presentation $\rho^{(N)} = c \prod_i B_{i,i+1}$ with $c > 0$ and
    commuting nearest-neighbor factors. The projector-limit convention after
    Definition~4.8 of \cite{Cirac2017MPDO_AnnPhys} identifies this with the
    exponential Gibbs form.
\end{definition}

\begin{theorem}[GSNNCH iff commuting form]
    \label{thm:mpdo_is_gsnnch_iff_has_commuting_form}
    \lean{MPOTensor.isGSNNCH_iff_hasCommutingForm}
    \leanok
    \uses{def:mpdo_is_gsnnch, def:mpdo_has_commuting_form}
    For an MPO tensor, the GSNNCH condition is equivalent to the existence of
    commuting-form data on every finite chain.
\end{theorem}

\begin{proof}\leanok
A positive normalization constant converts commuting-form data into a GSNNCH
witness, and every GSNNCH witness remembers its underlying commuting-form data.
\end{proof}

\begin{definition}[GSNNCH with zero correlation length]
    \label{def:mpdo_is_gsnnch_zcl}
    \lean{MPOTensor.IsGSNNCHWithZCL}
    \leanok
    \uses{def:mpdo_is_gsnnch, def:mpo_is_zcl}
    This is the conjunction of the GSNNCH condition with MPO zero correlation
    length, i.e. the content of item~(iii) in the simple-MPDO theorem.
\end{definition}

\begin{theorem}[GSNNCH with zero correlation length iff commuting form with zero correlation length]
    \label{thm:mpdo_is_gsnnch_zcl_iff_has_commuting_form_and_is_zcl}
    \lean{MPOTensor.isGSNNCHWithZCL_iff_hasCommutingForm_and_isZCL}
    \leanok
    \uses{def:mpdo_is_gsnnch_zcl, thm:mpdo_is_gsnnch_iff_has_commuting_form}
    The GSNNCH--ZCL branch is equivalent to the conjunction of the
    commuting-form property with MPO zero correlation length.
\end{theorem}

\begin{proof}\leanok
Unfold the definition of GSNNCH with zero correlation length and apply
Theorem~\ref{thm:mpdo_is_gsnnch_iff_has_commuting_form}.
\end{proof}

\begin{remark}[Remaining step toward Proposition~C.6]
The remaining entropy-side theorem is to derive the commuting-form property
from the strong-area-law hypotheses of Appendix~C.2. The present file isolates
its GSNNCH target once that step is available.
\end{remark}
